\chapter[Fundamentação Teórica]{Fundamentação Teórica}

Neste capítulo são apresentados os conceitos para o desenvolvimento deste trabalho, abrangendo aspectos do sistema cardiovascular, 
eletrocardiograma, processamento de sinais, inteligência artificial e, em especial, redes neurais convolucionais. A fundamentação 
teórica fornece a base necessária para compreender as etapas de modelagem e classificação de arritmias cardíacas em registros de ECG.

\section{Sistema cardiovascular e arritmias cardíacas}

O sistema cardiovascular é responsável pelo transporte de oxigênio, nutrientes e substâncias essenciais para o funcionamento do organismo, além de 
remover metabólitos através da circulação sanguínea. O coração atua como uma bomba eletromecânica, cujo funcionamento depende da geração e condução 
de impulsos elétricos iniciados no nó sinoatrial e propagados por estruturas especializadas do miocárdio \citeonline{Guyton}.

Em condições normais, a condução elétrica ocorre de forma organizada, resultando em um ritmo cardíaco regular, denominado ritmo sinusal. Alterações 
nesse processo caracterizam as arritmias cardíacas, que são distúrbios na formação ou na condução do impulso elétrico, podendo gerar batimentos rápidos, 
lentos ou irregulares. Essas alterações podem ter origem atrial, ventricular ou nodal, e variam desde eventos benignos até condições graves, como 
taquicardia ventricular e fibrilação ventricular, associadas a elevado risco de mortalidade \citeonline{Samesima}.

\section{Eletrocardiograma}

O ECG é um exame não invasivo amplamente utilizado para o monitoramento da atividade elétrica cardíaca, registrando as variações de potencial elétrico 
geradas durante o ciclo cardíaco. O sinal obtido é representado graficamente por um traçado composto pelas ondas P, QRS e T, bem como pelos intervalos 
PR e QT, que fornecem informações relevantes sobre o ritmo cardíaco, a condução elétrica e possíveis alterações funcionais e estruturais do coração.

\begin{figure}[ht]
  \caption{Eletrocardiograma normal}
  \label{fig: ECGnormal}
  \centering
  \includegraphics[width=0.7\textwidth]{ECG.png}
  \legend{\citeonline{Guyton}}
  
\end{figure}

Na figura \ref{fig: ECGnormal} é possível ver um ECG normal, onde é identificado o complexo QRS, as ondas P e T, os intervalos PR e QT, e o segmento ST. 
Apesar da grande utilização desse exame para diagnósticos, ele é sujeito a ruídos, como interferência da rede elétrica, artefatos musculares, variação da 
linha de base, variação entre pacientes. O que acaba impactando tanto na análise manual quanto na análise automatizada desse exame.

\section{Inteligência artificial}

A Inteligência Artificial (IA) refere-se ao desenvolvimento e aplicação de algoritmos capazes de identificar e aprender padrões complexos a partir de 
grandes volumes de dados. Essas técnicas permitem que sistemas computacionais executem tarefas cognitivas, como previsão, classificação e detecção de 
anomalias, frequentemente alcançando desempenho comparável ou superior ao de especialistas humanos \citeonline{Jiang}.

Com o avanço da tecnologia, houve um aumento de dispositivos de monitoramento de saúde, notadamente os tecnologias vestíveis (wearables), como 
relógios inteligentes e pulseiras de atividade. Esses dispositivos possibilitam o registro contínuo e não invasivo de sinais fisiológicos, gerando um 
volume massivo de dados (Big Data) que ultrapassa a capacidade humana de análise manual em tempo hábil \citeonline{Majumder2017}. Com esse cenário, 
o desenvolvimento de modelos automatizados robustos, capazes de processar esses sinais para detectar e classificar arritmias cardíacas se fazem necessários.

\section{Redes neurais convolucionais}

As Redes Neurais Convolucionais (CNNs) representam uma classe especializada de redes neurais artificiais projetadas especificamente para processar 
dados que possuem uma topologia de grade conhecida, como séries temporais (1D), no caso de sinais de ECG, ou dados de imagem (2D), como 
espectrogramas \citeonline{Goodfellow2016}.

Diferentemente das redes neurais tradicionais (Perceptrons Multicamadas), onde cada neurônio de uma camada está conectado a todos os neurônios da 
camada seguinte, as CNNs utilizam a operação matemática de convolução no lugar da multiplicação matricial convencional em pelo menos uma de suas 
camadas. Essa arquitetura permite que o modelo aprenda filtros (ou kernels) que deslizam sobre os dados de entrada, extraindo características locais 
e preservando a relação espacial ou temporal dos dados.

Matematicamente, em bibliotecas de aprendizado profundo, a operação de convolução unidimensional discreta aplicada a uma entrada $x$ com um 
filtro $w$ de tamanho $K$, pode ser expressa pela seguinte equação:

\begin{equation}
    \label{eq: eq1}
	y[i] = \sum_{k=0}^{K-1} x[i+k] \cdot w[k]
\end{equation}

Onde $x[i+k]$ representa o valor do sinal de entrada deslocado, $w[k]$ é o peso do filtro na posição $k$, e $y[i]$ é o valor resultante na saída. 
O conjunto desses valores de saída compõe o que é chamado de mapa de características (\textit{feature map}). O mapa de características atua como um 
ativador de padrões, destacando as regiões do sinal original onde as características buscadas pelo filtro, como a inclinação do complexo QRS em 
um batimento, estão mais presentes.

Após a operação de convolução, os mapas de características são invariavelmente submetidos a uma função de ativação não linear. No presente trabalho, 
adota-se a função Unidade Linear Retificada (ReLU), definida por:

\begin{equation}
    \label{eq: eq2}
    f(x) = \max(0, x)
\end{equation}

A aplicação da ReLU garante que valores negativos resultantes da convolução sejam descartados (zerados), enquanto valores positivos sejam mantidos 
inalterados. Essa introdução de não linearidade é estritamente necessária para que a rede consiga mapear funções complexas. Além disso, a ReLU é 
amplamente adotada por sua eficiência computacional, permitindo o treinamento eficaz de arquiteturas profundas \citeonline{Goodfellow2016}.

Uma arquitetura típica de CNN é composta por uma sequência de três tipos principais de camadas: camadas convolucionais, responsáveis pela extração de 
característica. Camadas de subamostragem (pooling), que reduzem a dimensionalidade dos mapas de características e conferem invariância a pequenas 
translações. E camadas totalmente conectadas (dense layers), que realizam a classificação final com base nos atributos extraídos. Essa estrutura 
hierárquica permite que a rede aprenda desde padrões simples, como bordas e picos, nas primeiras camadas, até conceitos abstratos e complexos nas 
camadas mais profundas \citeonline{Yamashita}.


\section{Preparação e Representação dos Dados}

O desempenho de algoritmos de aprendizado profundo é fortemente influenciado pela forma como os dados são representados e pré-processados. Em sinais 
biomédicos como o ECG, os dados brutos frequentemente apresentam variações de amplitude e ruídos de linha de base que podem dificultar a convergência 
do modelo durante o treinamento. Portanto, técnicas de normalização, como a padronização Z-score ou a normalização Min-Max, são fundamentais para 
garantir que as características de entrada estejam na mesma escala, facilitando a otimização dos pesos da rede.

Além do pré-processamento no domínio do tempo, a representação adequada do sinal é crucial para explorar diferentes arquiteturas de redes neurais. 
Enquanto as CNNs 1D operam diretamente sobre a série temporal, as CNNs 2D exigem uma transformação dos dados para o domínio tempo-frequência.

Para converter o sinal unidimensional em uma representação bidimensional (imagem), utiliza-se frequentemente a Transformada de Fourier de Curta 
Duração (STFT). O ECG é um sinal não estacionário, ou seja, suas frequências variam ao longo do tempo. A Transformada de Fourier clássica captura 
apenas o conteúdo frequencial global, perdendo a informação de quando essas frequências ocorrem. A STFT supera essa limitação ao dividir o sinal 
em janelas temporais curtas e aplicar a transformada em cada segmento, gerando um espectrograma que representa a densidade espectral de energia ao 
longo do tempo \citeonline{Stankovic2016}. Essa representação visual permite que a CNN2D identifique padrões complexos, como a variação abrupta de 
frequência no complexo QRS.

A STFT para um sinal de tempo discreto $x[n]$ é definida pela equação:

\begin{equation}
    \label{eq: eq3}
    X(m, \omega) = \sum_{n=-\infty}^{\infty} x[n] \cdot w[n-m] \cdot e^{-j\omega n}
\end{equation}

Onde $w[n]$ representa a função de janela centrada no instante de tempo $m$, e $\omega$ é a frequência angular. O quadrado da 
magnitude dessa transformada, $|X(m, \omega)|^2$, gera o espectrograma, que representa a densidade espectral de energia ao longo do 
tempo.

Um aspecto crítico na aplicação da STFT é o compromisso fundamental entre a resolução temporal e a resolução frequencial, governado pelo Princípio da 
Incerteza de Heisenberg (conhecido em processamento de sinais como limite de Gabor). Esse princípio estabelece que é matematicamente impossível obter 
precisão arbitrária simultaneamente no tempo e na frequência.

A escolha do tamanho da janela $w[n]$ dita esse balanço: uma janela estreita proporciona excelente localização no tempo (ideal para detectar o instante 
exato de ocorrência de uma anomalia), mas resulta em baixa resolução em frequência. Inversamente, uma janela larga melhora a precisão frequencial, 
mas "borra" o sinal no tempo. A definição dos parâmetros da STFT no pré-processamento do ECG é, portanto, um parâmetro de projeto essencial para 
equilibrar a visualização de variações abruptas de alta frequência (como o complexo QRS) e de oscilações mais lentas (como as ondas P e T).


\section{Divisão dos Dados e Validação}

Um dos desafios centrais no treinamento de redes neurais artificiais é o compromisso entre a otimização e a generalização. A generalização é a
capacidade do modelo de performar corretamente em dados inéditos. Quando um modelo se ajusta excessivamente aos dados de treino, capturando ruídos 
ou flutuações aleatórias em vez dos padrões subjacentes, ocorre o fenômeno conhecido como sobreajuste ou \textit{overfitting}.

Para diminuir esse problema e avaliar a capacidade de generalização de forma imparcial, a literatura estabelece a necessidade de particionar o conjunto 
de dados original. Uma das abordagens mais difundidas é o método \textit{Hold-out}. Nesta técnica, os dados são divididos aleatoriamente em subconjuntos 
distintos e independentes: o conjunto de treinamento, utilizado para o ajuste iterativo dos pesos da rede (aprendizado), e o conjunto de teste, que 
permanece isolado durante todo o processo de otimização, sendo utilizado exclusivamente para a avaliação final do desempenho do classificador.

\section{Regularização}

A regularização consiste em um conjunto de estratégias aplicadas aos algoritmos de aprendizado de máquina com o objetivo de reduzir o erro de 
generalização, mesmo que isso não reduza (ou até aumente ligeiramente) o erro de treinamento. Essas técnicas são fundamentais para prevenir o 
\textit{overfitting}, especialmente em redes neurais profundas com grande número de parâmetros \citeonline{Goodfellow2016}.

Uma das técnicas de regularização mais eficazes e amplamente utilizadas é o Dropout. Essa técnica consiste na desativação temporária e aleatória de 
uma porcentagem de neurônios (e suas conexões) durante a fase de treinamento. Isso força a rede a aprender características mais robustas e distribuídas, 
impedindo que os neurônios desenvolvam uma co-adaptação complexa onde um depende excessivamente da correção de erros de outros vizinhos. O Dropout aplica 
uma máscara binária $r$ aos neurônios de uma camada, onde cada elemento de $r$ segue uma distribuição de Bernoulli com probabilidade $p$:

\begin{align}
    \label{eq: eq4}
    r_j &\sim \text{Bernoulli}(p) \\
    \tilde{y} &= r \odot y \nonumber
\end{align}

Onde $y$ é o vetor de ativações originais e $\tilde{y}$ é o vetor resultante após o descarte. Essa operação força a rede a aprender características 
redundantes e distribuídas, impedindo a co-adaptação excessiva entre neurônios. No presente trabalho, utiliza-se uma taxa de $p = 0,5$ nas 
camadas densas. Esse valor é amplamente adotado na literatura, conforme proposto por \citeonline{Srivastava}, por apresentar bom desempenho empírico na 
redução de overfitting.

Do ponto de vista teórico, a variância da distribuição Bernoulli associada ao dropout, dada por $\mathrm{Var}(m) = p(1 - p)$, é maximizada 
em $p = 0{,}5$, o que implica maior aleatoriedade na ativação dos neurônios. No entanto, a escolha dessa taxa também depende do contexto do modelo e 
dos dados, sendo frequentemente definida com base em validação empírica.

Além do Dropout, utiliza-se a Batch Normalization (Normalização em Lote) para estabilizar o treinamento. Seu objetivo primário é mitigar o 
problema da mudança de covariável interna (\textit{internal covariate shift}), garantindo que as entradas de cada camada tenham uma distribuição estável. 
Para um minilote $\mathcal{B} = \{x_1, ..., x_m\}$, a operação é definida pelas seguintes etapas:

1. Cálculo da média do lote: \quad $\displaystyle \mu_{\mathcal{B}} = \frac{1}{m} \sum_{i=1}^{m} x_i$ \hfill \refstepcounter{equation}(\theequation)\label{eq:eq5}

2. Cálculo da variância do lote: \quad $\displaystyle \sigma_{\mathcal{B}}^2 = \frac{1}{m} \sum_{i=1}^{m} (x_i - \mu_{\mathcal{B}})^2$ \hfill \refstepcounter{equation}(\theequation)\label{eq:eq6}

3. Normalização: \quad $\displaystyle \hat{x}_i = \frac{x_i - \mu_{\mathcal{B}}}{\sqrt{\sigma_{\mathcal{B}}^2 + \epsilon}}$ \hfill \refstepcounter{equation}(\theequation)\label{eq:eq7}

4. Transformação afim: \quad $\displaystyle y_i = \gamma \hat{x}_i + \beta$ \hfill \refstepcounter{equation}(\theequation)\label{eq:eq8}

Onde $\epsilon$ é uma constante pequena para estabilidade numérica, e $\gamma$ (escala) e $\beta$ (deslocamento) são parâmetros aprendíveis durante o 
treinamento \citeonline{Ioffe}.

\section{Funções de Perda e Otimizadores}

O treinamento de uma rede neural é um problema de otimização cujo objetivo é minimizar uma função de custo (ou função de perda) que quantifica a 
disparidade entre as previsões do modelo e os rótulos reais dos dados. Para problemas de classificação multiclasse, onde as saídas representam uma 
distribuição de probabilidade sobre $C$ classes mutuamente exclusivas, a função de perda mais adequada e amplamente utilizada é a Entropia Cruzada 
Categórica (Categorical Cross-Entropy) \citeonline{Goodfellow2016}.

Matematicamente, essa função penaliza o modelo de forma logarítmica quando ele atribui uma probabilidade baixa à classe correta. Ao minimizar a 
entropia cruzada, a rede busca aproximar a distribuição de probabilidade predita (geralmente obtida via função de ativação Softmax na camada de saída) 
da distribuição real dos dados, maximizando assim a classificação correta \citeonline{Bishop2006}.

Para realizar a minimização dessa função de perda, utilizam-se algoritmos de otimização que ajustam iterativamente os pesos da rede com base nos 
gradientes calculados pelo algoritmo de Backpropagation. Embora o Gradiente Descendente Estocástico seja a base desses métodos, algoritmos com taxas 
de aprendizado adaptativas têm se mostrado mais eficientes em cenários complexos.

O otimizador Adam (Adaptive Moment Estimation), combina as propriedades do AdaGrad, que performa bem com gradientes esparsos, e do RMSProp, que lida 
com objetivos não estacionários. O algoritmo calcula taxas de aprendizado adaptativas individuais para cada parâmetro a partir das estimativas do 
primeiro e segundo momentos dos gradientes. O grande diferencial do Adam em relação aos seus 
predecessores é a inclusão de uma correção de viés (bias correction), que compensa o fato de os momentos serem inicializados como zero. Essa 
característica torna o Adam mais robusto nas etapas iniciais do treinamento e reduz a necessidade de ajuste fino (tuning) dos hiperparâmetros, 
garantindo uma convergência estável mesmo em cenários ruidosos \citeonline{Ruder2016}.

\section{Métricas de Avaliação}

A avaliação quantitativa de modelos de classificação é realizada através da análise da matriz de confusão. Esta matriz confronta as classes reais dos 
dados com as classes preditas pelo modelo, permitindo a extração de quatro valores fundamentais para cada classe: Verdadeiros Positivos ($TP$), 
Verdadeiros Negativos ($TN$), Falsos Positivos ($FP$) e Falsos Negativos ($FN$).

Os componentes fundamentais extraídos da matriz de confusão são definidos da seguinte forma, considerando a classe de interesse (por exemplo, 
a presença de uma arritmia) como a classe positiva:

\begin{itemize}
    \item \textbf{Verdadeiro Positivo (TP -- \textit{True Positive}):} Representa o acerto positivo do classificador. Ocorre quando o modelo prediz 
    corretamente a presença da condição de interesse (a arritmia existe e foi detectada).
    
    \item \textbf{Verdadeiro Negativo (TN -- \textit{True Negative}):} Representa o acerto negativo. Ocorre quando o modelo classifica corretamente a 
    ausência da condição (o batimento é normal e foi classificado como normal).
    
    \item \textbf{Falso Positivo (FP -- \textit{False Positive}):} Também conhecido como \textbf{Erro do Tipo I} ou "alarme falso". Ocorre quando o 
    modelo prediz incorretamente a presença da condição em uma amostra que é, na realidade, negativa (ex: classificar um batimento normal como arritmia).
    
    \item \textbf{Falso Negativo (FN -- \textit{False Negative}):} Também conhecido como \textbf{Erro do Tipo II}. Ocorre quando o modelo falha em 
    detectar a condição existente, classificando uma amostra positiva como negativa (ex: classificar uma arritmia perigosa como um batimento normal).
\end{itemize}

A partir desses valores, calculam-se métricas específicas que expressam diferentes aspectos do desempenho do classificador. A Acurácia é a métrica mais 
intuitiva, representando a proporção global de acertos do modelo sobre o total de amostras. Embora útil, ela pode ser enganosa em bases de dados 
desbalanceadas, onde o modelo pode obter alta acurácia apenas classificando todas as amostras na classe majoritária.

\begin{align}
    \label{eq: eq9}
    \text{Acurácia} = \frac{TP + TN}{TP + TN + FP + FN}
\end{align}

Em contextos médicos, a Sensibilidade (ou Recall) e a Especificidade assumem papel crucial. A Sensibilidade mede a capacidade do modelo de detectar 
corretamente os casos positivos (presença de arritmia), evitando falsos negativos que poderiam colocar o paciente em risco. Já a Especificidade avalia 
a capacidade de identificar corretamente os casos negativos (batimentos normais), evitando alarmes falsos.

\begin{align}
    \label{eq: eq10}
    \text{Sensibilidade} = \frac{TP}{TP + FN} \quad ; \quad \text{Especificidade} = \frac{TN}{TN + FP}
\end{align}

A Precisão avalia a confiabilidade das predições positivas, indicando a proporção de classificações corretas entre todas as amostras que o modelo 
classificou como positivas. Por fim, o F1-Score é a média harmônica entre a Precisão e a Sensibilidade. Esta métrica é particularmente valiosa para 
comparar modelos, pois penaliza valores extremos: um F1-Score alto só é obtido se tanto a Precisão quanto a Sensibilidade forem altas, garantindo um 
equilíbrio entre a detecção de arritmias e a minimização de falsos positivos.

\begin{align}
    \label{eq: eq11}
    \text{Precisão} = \frac{TP}{TP + FP} \quad ; \quad \text{F1-Score} = 2 \cdot \frac{\text{Precisão} \cdot \text{Sensibilidade}}{\text{Precisão} + \text{Sensibilidade}}
\end{align}